\subsection{Morphemes, Words and Compounds}
\label{sec:morphology}
\slidesrc{Terminology introduction: 35--38, 42}

\subsubsection{Morphemes}
\slidesrc{Terminology introduction: 35--36}

\begin{nbdef}{Morpheme}
A \textbf{morpheme} is the smallest grammatical unit in a language; in other words, the
smallest \emph{meaningful} unit of a language. The field of linguistics that studies
morphemes is called \textbf{morphology}.
\end{nbdef}

A morpheme is not the same as a word. The principal difference is that a morpheme may
or may not stand alone, whereas a word, by definition, is freestanding. In
\emph{dogs} there is one word but two morphemes: \emph{dog} (the animal) and \emph{-s}
(``more than one''). Neither \emph{d} nor \emph{og} carries meaning, so they are not
morphemes.

\begin{nbdef}{Free and bound morphemes}
Every morpheme is either free or bound.
\begin{itemize}[nosep]
  \item A \textbf{free morpheme} can function independently as a word (e.g.\
        \emph{town}, \emph{dog}).
  \item A \textbf{bound morpheme} appears only as part of a word, always together with a
        \textbf{root} and sometimes with other bound morphemes. For example,
        \emph{un-} never appears alone. Most bound morphemes in English are
        \textbf{affixes}: \textbf{prefixes} (before the root, \emph{un-}) and
        \textbf{suffixes} (after the root, \emph{-ness}).
\end{itemize}
\end{nbdef}

\begin{nbexample}{unbreakable}
\emph{Unbreakable} is made of three morphemes:
\begin{center}
\begin{tikzpicture}[
    m/.style={draw, rounded corners, minimum height=0.8cm, minimum width=1.7cm, font=\small},
    lab/.style={font=\scriptsize, align=center}]
  \node[font=\large\itshape] (w) at (2.4,1.5) {unbreakable};
  \node[m, fill=nbred!8,  draw=nbred]  (a) at (0,0)   {un-};
  \node[m, fill=nbblue!8, draw=nbblue] (b) at (2.4,0) {break};
  \node[m, fill=nbred!8,  draw=nbred]  (c) at (4.8,0) {-able};
  \draw[nbgray] (w.south) -- (a.north); \draw[nbgray] (w.south) -- (b.north); \draw[nbgray] (w.south) -- (c.north);
  \node[lab, below=2pt of a] {bound prefix\\``not''};
  \node[lab, below=2pt of b] {root\\free morpheme};
  \node[lab, below=2pt of c] {bound suffix\\``can be done''};
\end{tikzpicture}
\end{center}
\end{nbexample}

\begin{nbwarning}[\emph{-able} is bound]
The slide calls \emph{-able} ``a free morpheme signifying \emph{can be done}''. As a
suffix it is \textbf{bound}: it cannot stand alone with the meaning ``that can be
done'' and needs a root (\emph{break-able}, \emph{read-able}). The English word
\emph{able} exists, but it is a separate free morpheme with a different use (``she is
able to swim''). So \emph{unbreakable} = bound prefix + free root + bound suffix.
\end{nbwarning}

\subsubsection{Derivational and inflectional morphemes}
\slidesrc{Terminology introduction: 37}

Bound morphemes are further divided according to \emph{what they change}.

\begin{nbdef}{Derivational and inflectional morphemes}
\begin{itemize}[nosep]
  \item A \textbf{derivational morpheme}, combined with a root, changes the
        \textbf{meaning} or the \textbf{part of speech} of the word: it creates a
        \emph{new word}.
  \item An \textbf{inflectional morpheme} modifies a verb's tense, aspect, mood, person or
        number, or a noun's, pronoun's or adjective's number, gender or case,
        \emph{without} changing the word's meaning or class (part of speech): it creates
        a different \emph{form of the same word}.
\end{itemize}
\end{nbdef}

\begin{center}
\small
\begin{tabularx}{\linewidth}{@{}l l l X@{}}
\toprule
\textbf{Root} & \textbf{Morpheme} & \textbf{Result} & \textbf{What changed} \\
\midrule
\multicolumn{4}{@{}l}{\emph{Derivational}} \\
happy (adjective) & -ness & happiness (noun) & part of speech: adjective $\to$ noun \\
kind              & un-   & unkind           & meaning is inverted \\
\midrule
\multicolumn{4}{@{}l}{\emph{Inflectional}} \\
dog  & -s  & dogs   & number (singular $\to$ plural) \\
wait & -ed & waited & tense (present $\to$ past) \\
\bottomrule
\end{tabularx}
\captionof{table}{Derivation makes a new word, inflection a new form of the same word.}
\end{center}

\begin{nbintuition}[a dictionary test]
Ask: ``would a dictionary give the result its own entry?'' \emph{Happiness} and
\emph{unkind} have their own entries (derivation), \emph{dogs} and \emph{waited} are found
under \emph{dog} and \emph{wait} (inflection). This is exactly the difference between
\emph{lemmatisation} (going back to the dictionary form, removing inflection) and the
analysis of word formation, which we will meet later in the course.
\end{nbintuition}

\subsubsection{Words}
\slidesrc{Terminology introduction: 38}

\begin{nbdef}{Word}
A \textbf{word} is the smallest element that can be uttered in isolation with
objective or practical meaning.
\end{nbdef}

This contrasts with a morpheme, which is the smallest unit of meaning but does not
necessarily stand on its own:
\begin{itemize}[nosep]
  \item a word may consist of a \textbf{single morpheme}: \emph{oh!}, \emph{rock},
        \emph{red}, \emph{quick}, \emph{run}, \emph{expect};
  \item or of \textbf{several}: \emph{rock-s}, \emph{red-ness}, \emph{quick-ly},
        \emph{runn-ing}, \emph{un-expect-ed};
  \item and a morpheme may be unable to stand on its own as a word: in the examples,
        \emph{-s}, \emph{-ness}, \emph{-ly}, \emph{-ing}, \emph{un-}, \emph{-ed}.
\end{itemize}

\subsubsection{Compounds}
\slidesrc{Terminology introduction: 42}

\begin{nbdef}{Compound}
A \textbf{compound} is a word composed of more than one \emph{free} morpheme.
\end{nbdef}

The last part is usually the \textbf{head} (it gives the category of the whole word),
and the first part is the \textbf{modifier}. Compounds can be written as one word
(\emph{football}), with a hyphen (\emph{blue-green}) or even as two words
(\emph{snow white}).

\begin{center}
\small
\begin{tabular}{@{}l l l @{\hspace{1.6cm}} l l l@{}}
\toprule
\textbf{Modifier} & \textbf{Head} & \textbf{Compound} & \textbf{Modifier} & \textbf{Head} & \textbf{Compound} \\
\midrule
noun        & noun      & football    & noun        & verb        & browbeat \\
adjective   & noun      & blackboard  & adjective   & verb        & highlight \\
verb        & noun      & breakwater  & verb        & verb        & freeze-dry \\
preposition & noun      & underworld  & preposition & verb        & undercut \\
noun        & adjective & snow white  & noun        & preposition & love-in \\
adjective   & adjective & blue-green  & adverb      & preposition & forthwith \\
verb        & adjective & tumbledown  & verb        & preposition & takeout \\
preposition & adjective & over-ripe   & preposition & preposition & without \\
\bottomrule
\end{tabular}
\captionof{table}{English compounds by the category of their two parts.}
\end{center}

\begin{nbwarning}[compound vs.\ derived word]
\emph{Blackboard} is a compound (\emph{black} and \emph{board} are both free).
\emph{Unkind} is \emph{not} a compound: \emph{un-} is bound, so it is a derived word.
Also note that the meaning of a compound is often not the sum of its parts: a
\emph{blackboard} can be green.
\end{nbwarning}
